\chapter{Blepharitis}
\index{Blepharitis}
\label{sec:Blepharitis}

\section{Overview}
Blepharitis is a chronic inflammatory condition affecting the eyelid margins, characterized by redness, crusting, and irritation. It is one of the most common ophthalmic conditions encountered in clinical practice, affecting people of all ages. The condition is often associated with meibomian gland dysfunction, seborrheic dermatitis, and rosacea.
\section{Classification}

\begin{description}[font=\normalfont\bfseries,leftmargin=3em,labelwidth=2.5em]
  \item[Cause] Inflammatory / Infectious
  \item[Organ System] Eyes
  \item[Pathophysiology] Meibomian gland dysfunction with bacterial overgrowth and inflammation of eyelid margins
  \item[Duration] Chronic (months to years)
  \item[Onset] Gradual
  \item[Mode of Transmission] Non-communicable
  \item[Clinical Course] Chronic, relapsing-remitting
  \item[Severity] Mild to Moderate
  \item[Extent of Disease] Localized
  \item[Age Group] All ages
  \item[Epidemiology] Very common; estimated prevalence 37--47\% in ophthalmology patients
  \item[ICD-11 Code] 9A61.0
\end{description}

\section{Causes}
Anterior blepharitis involves inflammation at the base of the eyelashes, often from bacterial colonization (Staphylococcus aureus, S. epidermidis) or seborrheic dermatitis. Posterior blepharitis results from meibomian gland dysfunction with altered lipid secretion, creating an environment that promotes bacterial growth and inflammation. Demodex mite infestation contributes in some cases.

\section{Risk Factors}

\begin{itemize}[noitemsep]
  \item Seborrheic dermatitis
  \item Rosacea
  \item Meibomian gland dysfunction
  \item Contact lens wear
  \item Dry eye syndrome
  \item Advanced age
  \item Poor eyelid hygiene
  \item Diabetes mellitus
\end{itemize}

\section{Signs and Symptoms}

\subsection{Early symptoms}
\begin{itemize}[noitemsep]
  \item Early manifestations of blepharitis may be subtle and nonspecific
\end{itemize}

\subsection{Common symptoms}
\begin{itemize}[noitemsep]
  \item \subsection{Early symptoms}
  \item \subsection{Common symptoms}
  \item \textbf{Common symptoms:}
  \item Chronic redness and swelling of the eyelid margins
  \item Crusting, flaking, and collarettes at the base of eyelashes
  \item Itching, burning, and foreign body sensation
  \item Tearing and dry eye symptoms
  \item Photophobia and blurred vision
  \item Misdirected or loss of eyelashes
  \item Pain or discomfort in the affected area
  \item Difficulty performing daily activities
\end{itemize}

\subsection{Less common symptoms}
\begin{itemize}[noitemsep]
  \item Atypical or less common presentations of blepharitis may occur
\end{itemize}

\subsection{Emergency warning signs}
\begin{itemize}[noitemsep]
  \item Severe or worsening symptoms of blepharitis require immediate medical evaluation
\end{itemize}
\section{Progression and Stages}
Blepharitis follows a chronic relapsing-remitting course with exacerbations triggered by stress, poor hygiene, or underlying skin conditions. Without consistent management, inflammation can lead to structural changes in the eyelid margin, including scarring, lash loss (madarosis), and trichiasis.

\section{Complications}

\begin{itemize}[noitemsep]
  \item Recurrent hordeola and chalazia
  \item Dry eye syndrome from tear film disruption
  \item Conjunctival injection and keratitis
  \item Corneal ulceration and scarring (rare)
  \item Trichiasis (inward-turning lashes)
  \item Reduced quality of life
  \item Emotional and psychological effects
\end{itemize}

\section{Diagnosis}

\begin{itemize}[noitemsep]
  \item Clinical diagnosis by slit-lamp examination of eyelid margins
  \item Assessment of meibomian gland expression and secretion quality
  \item Lid margin swab for culture in refractory cases
  \item Demodex mite evaluation by lash epilation and microscopy
  \item Lab tests to check for underlying conditions
\end{itemize}

\section{Prevention}

\begin{itemize}[noitemsep]
  \item Daily eyelid hygiene with warm compresses and lid scrubs
  \item Manage underlying skin conditions (seborrhea, rosacea)
  \item Regular eyelid massage to express meibomian glands
  \item Avoid eye makeup contamination and replace products regularly
  \item Go for regular check-ups so any problems are caught early
\end{itemize}

\section{Treatment Options}

\begin{itemize}[noitemsep]
  \item Warm compresses applied daily for 5--10 minutes followed by lid massage
  \item Eyelid margin scrubs with diluted baby shampoo or commercial lid wipes
  \item Topical antibiotic ointments (erythromycin, bacitracin) for anterior blepharitis
  \item Topical cyclosporine or corticosteroids for inflammatory component
  \item Oral doxycycline or azithromycin for meibomian gland dysfunction
  \item Omega-3 fatty acid supplementation
  \item Artificial tears for associated dry eye
  \item Lifestyle changes and self-care
\end{itemize}

\section{Living with It}

\begin{itemize}[noitemsep]
  \item Follow your treatment plan as prescribed
  \item Keep up with regular appointments with your healthcare provider
  \item Adjust your daily habits as needed
  \item Reach out to family, friends, or support groups for help
  \item Keep learning about your condition and how to manage it
\end{itemize}

\section{Outlook}
Blepharitis is a chronic condition that typically requires long-term management rather than cure. With consistent eyelid hygiene and appropriate treatment, most patients achieve adequate symptom control. The condition tends to recur and requires ongoing maintenance therapy to prevent complications and preserve ocular surface health.
